\documentclass[11pt,a4paper]{article}
\usepackage[italian]{babel}
\usepackage[utf8]{inputenc}
\usepackage[T1]{fontenc}
\usepackage{amsmath,amssymb,amsthm}
\usepackage{geometry}
\geometry{margin=2.2cm}
\usepackage{enumitem}
\usepackage{tikz}
\usepackage{tabularx,array,booktabs}
\usepackage{xcolor}
\usepackage{hyperref}
\hypersetup{colorlinks=true, linkcolor=blue, urlcolor=blue}
\DeclareMathOperator{\arccot}{arccot}
\newcommand{\risp}{\vspace{2pt}\noindent\textbf{Soluzione.}\ }
\newcolumntype{Y}{>{\raggedright\arraybackslash}X}
\newcommand{\griglia}[1]{%
\begin{tikzpicture}[scale=0.75]
\draw[step=0.5,gray!30,very thin] (-4,-2.5) grid (4,4);
\draw[->] (-4,0) -- (4.2,0) node[right]{\small $x$};
\draw[->] (0,-2.5) -- (0,4.2) node[above]{\small $y$};
\node[below left] at (0,0) {\scriptsize $O$};
\node at (-2.4,3.6) {\small #1};
\end{tikzpicture}}
\newcommand{\vuoto}{\hspace{1.4cm}}

\begin{document}
\begin{center}
{\Large\bfseries Le funzioni goniometriche inverse}\\[3pt]
Scheda di lavoro -- Classe quarta \qquad Liceo Scientifico ``G. Galilei''\\[6pt]
Nome e cognome \rule{6cm}{0.4pt} \qquad Data \rule{3cm}{0.4pt}
\end{center}

\noindent\textbf{Ricorda.} Una funzione \`e invertibile se \`e biiettiva. Il grafico di $f^{-1}$ si ottiene da quello di $f$ scambiando le coordinate di ogni punto, cio\`e con la simmetria rispetto alla bisettrice $y=x$. Usa sempre \textbf{la stessa unit\`a di misura sui due assi} (suggerimento: $1=2$ quadretti, quindi $\tfrac{\pi}{2}\approx 3$ quadretti).

\section*{Parte A -- Costruiamo i grafici (in classe)}

\subsection*{A.1 Arcoseno}
Restringiamo il seno all'intervallo $\left[-\tfrac{\pi}{2},\tfrac{\pi}{2}\right]$. Completa la tabella, poi scambia le coordinate.

\begin{center}
\renewcommand{\arraystretch}{1.5}
\begin{tabular}{|c|*{7}{c|}}
\hline
$x$ & $-\frac{\pi}{2}$ & $-\frac{\pi}{3}$ & $-\frac{\pi}{6}$ & $0$ & $\frac{\pi}{6}$ & $\frac{\pi}{3}$ & $\frac{\pi}{2}$\\\hline
$\sin x$ & \vuoto&\vuoto&\vuoto&\vuoto&\vuoto&\vuoto&\vuoto\\\hline
\end{tabular}
\end{center}
Scambiando le due righe ottieni i punti del grafico di $y=\arcsin x$. Disegna nello stesso riferimento $y=\sin x$ (ristretta), la bisettrice $y=x$ e $y=\arcsin x$.

\subsection*{A.2 Arcocoseno}
Perch\'e per il coseno \emph{non} va bene l'intervallo $\left[-\tfrac{\pi}{2},\tfrac{\pi}{2}\right]$?\\[2pt] \rule{\textwidth}{0.4pt}

Restringiamo il coseno a $[0,\pi]$:
\begin{center}
\renewcommand{\arraystretch}{1.5}
\begin{tabular}{|c|*{7}{c|}}
\hline
$x$ & $0$ & $\frac{\pi}{6}$ & $\frac{\pi}{3}$ & $\frac{\pi}{2}$ & $\frac{2\pi}{3}$ & $\frac{5\pi}{6}$ & $\pi$\\\hline
$\cos x$ & \vuoto&\vuoto&\vuoto&\vuoto&\vuoto&\vuoto&\vuoto\\\hline
\end{tabular}
\end{center}

\begin{center}
\griglia{A.1} \qquad \griglia{A.2}
\end{center}

\subsection*{A.3 Arcotangente e arcocotangente}
Restringiamo la tangente a $\left(-\tfrac{\pi}{2},\tfrac{\pi}{2}\right)$ e la cotangente a $(0,\pi)$.
\begin{center}
\renewcommand{\arraystretch}{1.5}
\begin{tabular}{|c|*{5}{c|}}
\hline
$x$ & $-\frac{\pi}{3}$ & $-\frac{\pi}{4}$ & $0$ & $\frac{\pi}{4}$ & $\frac{\pi}{3}$\\\hline
$\tan x$ & \vuoto&\vuoto&\vuoto&\vuoto&\vuoto\\\hline
\end{tabular}
\qquad
\begin{tabular}{|c|*{5}{c|}}
\hline
$x$ & $\frac{\pi}{6}$ & $\frac{\pi}{4}$ & $\frac{\pi}{2}$ & $\frac{3\pi}{4}$ & $\frac{5\pi}{6}$\\\hline
$\cot x$ & \vuoto&\vuoto&\vuoto&\vuoto&\vuoto\\\hline
\end{tabular}
\end{center}
Che cosa diventano gli asintoti verticali $x=\pm\tfrac{\pi}{2}$ della tangente nel grafico dell'arcotangente? \rule{4.5cm}{0.4pt}

\begin{center}
\griglia{$\arctan$} \qquad \griglia{$\operatorname{arccot}$}
\end{center}

\subsection*{A.4 Tabella riassuntiva}
\begin{center}
\renewcommand{\arraystretch}{1.9}
\begin{tabular}{|l|p{2.6cm}|p{2.6cm}|p{2.6cm}|p{2.6cm}|}
\hline
 & $y=\arcsin x$ & $y=\arccos x$ & $y=\arctan x$ & $y=\arccot x$\\\hline
Dominio & & & & \\\hline
Codominio & & & & \\\hline
Monotonia & & & & \\\hline
Simmetrie & & & & \\\hline
Asintoti & & & & \\\hline
\end{tabular}
\end{center}

\section*{Parte B -- Valori e composizioni (in classe, a coppie)}
\begin{enumerate}[label=\textbf{B.\arabic*}]
\item Calcola:
\begin{enumerate}[label=\alph*), itemsep=2pt]
\begin{minipage}[t]{0.42\textwidth}
\item $\arcsin\frac12$
\item $\arcsin\left(-\frac{\sqrt3}{2}\right)$
\item $\arccos\left(-\frac{\sqrt2}{2}\right)$
\item $\arccos 0$
\end{minipage}
\begin{minipage}[t]{0.42\textwidth}
\item $\arctan\left(-\sqrt3\right)$
\item $\arctan 1$
\item $\arccot(-1)$
\item $\arccot\sqrt3$
\end{minipage}
\end{enumerate}

\item Calcola, se possibile. \emph{Attenzione: il risultato di una funzione inversa deve stare nel suo codominio!}
\begin{enumerate}[label=\alph*), itemsep=2pt]
\begin{minipage}[t]{0.42\textwidth}
\item $\sin\left(\arcsin\frac13\right)$
\item $\arcsin\left(\sin\frac{5\pi}{6}\right)$
\item $\arccos\left(\cos\left(-\frac{\pi}{3}\right)\right)$
\item $\arctan\left(\tan\frac{3\pi}{4}\right)$
\end{minipage}
\begin{minipage}[t]{0.42\textwidth}
\item $\cos\left(\arcsin\frac35\right)$
\item $\tan\left(\arccos\left(-\frac12\right)\right)$
\item $\sin\left(\arcsin 2\right)$
\item[$\ast$h)] $\sin\left(\arctan 2\right)$
\end{minipage}
\end{enumerate}
\end{enumerate}

\section*{Parte C -- Domini (a casa)}
Determina il dominio delle seguenti funzioni.
\begin{enumerate}[label=\alph*), itemsep=3pt]
\begin{minipage}[t]{0.42\textwidth}
\item $y=\arcsin(2x-1)$
\item $y=\arccos\left(x^2-1\right)$
\item $y=\arctan\dfrac{1}{x-2}$
\end{minipage}
\begin{minipage}[t]{0.42\textwidth}
\item $y=\sqrt{\arcsin x}$
\item $y=\arccot\sqrt{x}$
\item[$\ast$f)] $y=\arcsin\dfrac{x+1}{x-1}$
\end{minipage}
\end{enumerate}

\section*{Parte D -- Grafici deducibili (a casa)}
Partendo dai grafici della parte A, disegna con le trasformazioni geometriche note e indica dominio e codominio:
\begin{enumerate}[label=\alph*), itemsep=2pt]
\begin{minipage}[t]{0.42\textwidth}
\item $y=\arcsin x+\frac{\pi}{2}$
\item $y=2\arctan x$
\end{minipage}
\begin{minipage}[t]{0.42\textwidth}
\item $y=\arccos(x-1)$
\item $y=\left|\arcsin x\right|$
\end{minipage}
\end{enumerate}
Confronta il grafico a) con quello di $y=\arccos x$ ribaltato rispetto a\dots\ Che cosa osservi?

\section*{Parte E -- Equazioni e un problema (a casa)}
\begin{enumerate}[label=\textbf{E.\arabic*}]
\item Risolvi, motivando i casi impossibili:
\begin{enumerate}[label=\alph*), itemsep=2pt]
\begin{minipage}[t]{0.42\textwidth}
\item $\arcsin x=\frac{\pi}{3}$
\item $\arccos(2x)=\frac{2\pi}{3}$
\item $\arctan(x-1)=-\frac{\pi}{4}$
\end{minipage}
\begin{minipage}[t]{0.42\textwidth}
\item $\arccos x=-\frac{\pi}{6}$
\item $\arcsin x=\arccos x$
\end{minipage}
\end{enumerate}

\item \textbf{Una rampa a norma.} La normativa italiana sull'abbattimento delle barriere architettoniche (D.M.~236/1989) fissa per le rampe una pendenza massima dell'$8\%$, cio\`e un dislivello di $8$ cm ogni metro di percorso \emph{orizzontale}.
\begin{enumerate}[label=\alph*), itemsep=2pt]
\item Quale funzione inversa fornisce l'angolo di inclinazione massimo? Calcolalo in radianti e in gradi (approssima al centesimo).
\item Una rampa \`e costituita da una pedana lunga $5$ m (misurata lungo il piano inclinato) che supera un dislivello di $40$ cm. L'inclinazione si calcola con la stessa funzione del punto a)? Stabilisci se la rampa \`e a norma.
\end{enumerate}

\item[$\ast$\textbf{E.3}] Dimostra che $\arcsin x+\arccos x=\dfrac{\pi}{2}$ per ogni $x\in[-1,1]$.\\
\emph{Suggerimento:} poni $\alpha=\arcsin x$ e verifica che $\frac{\pi}{2}-\alpha$ sta in $[0,\pi]$ e ha coseno uguale a $x$.
\end{enumerate}

\end{document}
